% 2_thesis_structure.tex
\section{Structure of the Dissertation and Contributions}
\label{sec:ch1_dissertation_structure}

% Overview of dissertation organization
This section provides a roadmap through the dissertation's six chapters, explaining how each contributes to the integrated intelligent orchestration framework. The structure reflects the logical progression from foundational concepts through three interdependent technical contributions to synthesis and future directions.

% Chapter 2: Background and state of the art
Chapter~\ref{ch:background} provides the theoretical foundation for the dissertation. Part I introduces the core concepts: time-series forecasting for workload prediction, transfer learning and attention mechanisms for cross-application generalization, and multi-agent architectures for autonomous orchestration. Part II analyzes the state of the art in each area, reviewing existing edge-deployable models, generalization approaches, and autonomous orchestration systems to identify the specific gaps addressed by this dissertation.

% Chapter 3: Edge-deployable workload prediction
Chapter~\ref{ch:adaptive-orchestration} addresses the computational intensity limiting edge deployment through \ac{AERO}. This lightweight forecasting model achieves accurate workload prediction with only 599 parameters, representing a 99.98\% reduction compared to the best-performing transformer baseline (Pathformer). This enables deployment on resource-constrained edge devices where traditional transformer-based models requiring 500K+ parameters would be computationally impractical. The core technique is adaptive period detection that dynamically identifies dominant periodicities in workload data. This approach uses the inherent cyclical nature of microservice demands to maintain prediction accuracy with minimal parameters. Validation on real-world datasets from Google Borg and Alibaba production systems shows comparable accuracy to state-of-the-art models. Beyond accuracy, the approach achieves sub-millisecond inference (0.38ms), enabling proactive orchestration decisions at the edge with sub-50ms latency. This directly addresses Technical Objective 1 (Edge-Deployable Workload Prediction).

% Chapter 4: Cross-application generalization
With edge-deployable prediction established, Chapter~\ref{ch:attention-driven-prediction} addresses the generalization problem that necessitates per-application customization through \ac{OmniFORE}. This framework enables single models to serve multiple heterogeneous application types without per-application training or customization. This avoids the substantial operational overhead of maintaining hundreds of specialized models. The approach combines temporal clustering to discover fundamental workload patterns that transcend application-specific implementations with attention mechanisms that learn which patterns apply to which services. Evaluation on diverse workload traces demonstrates cross-dataset generalization, achieving 30\% better \ac{MAE} than \ac{ModernTCN} on Alibaba traces after training exclusively on Google Borg data. The framework maintains inference speeds suitable for real-time orchestration. This enables scalable deployment across hundreds of microservices without retraining overhead, satisfying Technical Objective 2 (Cross-Application Generalization).

% Chapter 5: Safe autonomous orchestration
Given these prediction capabilities, Chapter~\ref{ch:distributed-agentic-ai} introduces \ac{AgentEdge}, a distributed multi-agent framework that enables safe autonomous orchestration through digital twin simulation. The framework introduces the \ac{PARES} capability framework defining minimum agent requirements and coordinates four specialized agents through structured workflows using the \ac{ActSimCrit} methodology that validates orchestration strategies through simulation before deployment. The chapter systematically evaluates five leading open-source \acp{LLM} for orchestration tasks and compares \ac{AgentEdge} against baseline agentic frameworks (\ac{ReAct}, \ac{LATS}). Evaluation across six orchestration scenarios spanning 8 to 35 node infrastructures demonstrates 78.3\% success rates, 2.76× improvement over baseline approaches, and 10× reduction in \ac{API} call variability. Energy optimization scenarios achieve power savings reaching 300.8W. This fulfills Technical Objective 3 (Safe Autonomous Orchestration).

% Chapter 6: Synthesis and future directions
Having presented the three technical contributions, Chapter~\ref{ch:conclusions} synthesizes the dissertation's achievements and establishes the foundation for future research. This concluding chapter validates the integrated approach where \ac{AERO} provides edge-deployable forecasting, \ac{OmniFORE} enables cross-application generalization, and \ac{AgentEdge} delivers autonomous orchestration. The chapter assesses whether stated technical objectives were achieved through quantitative evidence from the validation studies. Finally, the chapter identifies future research directions organized across three thematic areas: reliability and trust (evaluation benchmarks, explainability, security), technical infrastructure (context engineering, inference efficiency, tool discovery, time-series integration), and system architecture (model-agnostic design, multi-agent coordination at scale).
